\textbf{Story/thesis:}

Existing multi-agent KG construction systems (KARMA, CooperKGC, OneKE)
treat the knowledge graph as a flat output artifact which agents
extract, a controller merges, and the pipeline terminates. The resulting
graph has no organizational memory: no record of which agent produced
which fact, no domain structure governing updates, and no mechanism for
routing queries to the knowledge owners. We argue that a knowledge graph
should not be a passive store but a \textbf{governed memory} partitioned
into domain-owned subgraphs where every triple is admitted through an
explicit governance decision, every decision is auditable, and
downstream applications like QA route through domain experts rather than
flat retrieval.

\begin{enumerate}
\def\labelenumi{\arabic{enumi}.}
\item
  What is, in one sentence, the takeaway for this whole \emph{system-}
  We propose the Governed Knowledge Graph, a data structure in which
  every triple is admitted through an explicit ownership-and-review
  protocol, and show that governance shapes graph structure and enables
  principled downstream reasoning.
\item
  Why this methodology or architecture is necessary

  \begin{enumerate}
  \def\labelenumii{\alph{enumii}.}
  \item
    Agents vs. other approaches?-\\
    1. Specialization reduces error modes. Entity extraction, relation
    extraction, and evidence verification have different failure modes.
    Specializing agents lets you diagnose and fix each independently.\\
    2. Conflict resolution requires arbitration. When two agents
    disagree on a triple, you need structured deliberation, not a single
    model\textquotesingle s confidence score. Multi-agent voting
    produces a decision with a rationale.\\
    3. Domain ownership requires distributed expertise. A single model
    has no principled reason to route a triple about methodology to a
    "methods" domain. Domain expert agents instantiate the ownership
    function $\phi$. A single LLM baseline can\textquotesingle t do any of
    these.
  \item
    Domain-expertise for data modularity - Each domain $d \in D$ is a
    coherent subgraph- a region of the graph with a common conceptual
    scope, a set of owning entities, and a relation schema. The
    ownership function $\phi$ partitions the graph into these regions. This
    enables: (a) modular updates- only the owning domain validates a
    proposed new triple; (b) modular queries- downstream applications
    route to the relevant expert rather than doing flat retrieval; (c)
    modular auditing- every decision is traceable to a specific domain
    owner.
  \end{enumerate}
\item
  Self-organization
\item
  KARMA and their claim

  \begin{enumerate}
  \def\labelenumii{\alph{enumii}.}
  \item
    Does it make sense to run a comparison to this system vs. KARMA?
    Maybe not
  \end{enumerate}
\end{enumerate}

\textbf{Abstract}

Existing multi-agent knowledge graph construction systems, including
KARMA and CooperKGC, treat the resulting graph as a flat artifact: a
centralized controller merges agent extractions without domain
ownership, auditable governance, or any mechanism for routing downstream
queries to the knowledge owners. The graph carries no organizational
memory and no record of which agent produced which fact. We propose
MaKG, a multi-agent framework that treats a knowledge graph as governed
memory, a structure partitioned into domain-owned subgraphs where every
triple is admitted through an explicit ownership-and-review protocol and
every admission decision is auditable.

MaKG introduces three coordinated components.

\begin{itemize}
\item
  First, a nine-stage extraction pipeline whose DomainClassifier
  discovers entity and relation types directly from document content,
  enabling the system to handle arbitrary domains without predefined
  ontologies followed by blackboard-based deliberation using a
  seven-level voting scale and structured debate for conflicting
  hypotheses.

  \begin{itemize}
  \item
    Call this the ``creation layer'' (the same term Pranav used) -
    documents are segmented, entities are extracted, relations are
    proposed, evidence is linked, and candidate triples are verified
  \end{itemize}
\item
  Second, a domain-expert QA system that routes queries through an
  automatically constructed organization chart built from the same
  ownership structure used during extraction, supporting iterative
  active graph exploration and claim-level provenance tracking.

  \begin{itemize}
  \item
    Lets call this the ``governance layer'' where the candidate triples
    are routed according to the domain ownership and admitted \emph{only
    after a governance decision}
  \end{itemize}
\item
  Third, KGAFE, an atomic-fact evaluation framework that decomposes QA
  answers into independently verifiable claims and checks each against
  the KG through a three-tier cascade exact triple match, multi-hop path
  entailment, and semantic neighborhood entailment producing provably
  correct evaluation sets from graph structure without human annotation.

  \begin{itemize}
  \item
    Kind of - lets propose this more as an application layer for QA
    (what \#2 kinda is): downstream tasks like QA and incremental
    enrichment (self-organization, etc.) using the governed structure
    than retrieving from a flat graph (this is key!)
  \end{itemize}
\end{itemize}

On the SciERC benchmark, governed triage substantially improves triple
quality and reduces hallucination relative to a flat extraction baseline
using identical agents and schema. In downstream QA, domain-routed
retrieval outperforms both document RAG and flat graph retrieval on
faithfulness, hallucination control, and coverage. The central
architectural insight is that MaKG\textquotesingle s ownership structure
is shared across both subsystems: the same domain partition that governs
which facts enter the graph determines how queries are routed
downstream, making governance the connective tissue between knowledge
creation and knowledge use rather than a post-hoc filter on either.

\textbf{Try something like this (adding more detail up front and making
the claim much more clear and cleaner):}

We introduce the Governed Knowledge Graph, a knowledge representation in
which facts are not merely stored as triples but admitted through
explicit ownership, review, and provenance decisions. This
organizational layer enables the same graph structure to support
controlled construction, incremental update, audit, and domain-routed
reasoning.

\textbf{Faizan Abstract}

On the SciERC benchmark, governed construction improves triple quality
relative to a flat extraction baseline while preserving entity quality
and reducing unsupported triples. In downstream QA, domain-routed
systems improve coverage and groundedness over document retrieval and
flat graph-use baselines.

\textbf{$\to$ Proving Data is stored in a better way}

\includegraphics[width=3.54312in,height=1.92839in]{image1.png}

\textbf{Figure 1: Our framework parses documents into a governed
knowledge graph through multi-agent deliberation and answers user
queries through domain-expert collaboration over the resulting graph.}

\includegraphics[width=6.5in,height=3.625in]{image2.png}

\textbf{Figure 2: System architecture.} The framework operates in two
modes over a shared governed knowledge graph. In \textbf{Learning mode}
(left), five worker agents run a sequential extraction pipeline:
documents are segmented, a domain schema is derived, and entities,
relations, and evidence are extracted in turn. Uncertain items posted by
extraction agents (Entity Extractor, Relation Extractor, Evidence
Linker) are routed through a deliberation side-loop where peer agents
vote and, if needed, debate before the items rejoin the main flow. Three
coordinator agents produce verified triples that are \emph{proposed} to
a governance board, which routes each triple as auto-approve, LLM
review, or escalate, and assigns an outcome (approve, revise, or reject)
under the $\gamma$ governance function before \emph{commit} to the graph. In
\textbf{Ask mode} (right), a QA orchestrator routes user queries to
domain experts that \emph{consult} the graph through active exploration,
with disagreements resolved by a debate arena and a critic agent before
a grounded answer is returned with provenance. SharedMemory and a typed
MessageBus underlie all agent communication.

\textbf{Governed Knowledge Graph}

Definition 1.\\
A Governed Knowledge Graph is a tuple $(G = (E, T, D, \phi, \gamma))$, where:

E is a set of entities.\\
$T \subseteq E \times R \times E$ is a set of typed triples.\\
$D = \{d_1, \ldots, d_k\}$ is a set of domains.\\
$\phi: E \to 2^D$ is an ownership function assigning each entity to one or
more governing domains.\\
γ: T → \{approve, reject, revise, escalate, auto\_approve\} is a
governance function over proposed triples.

This definition differs from a standard knowledge graph in two ways.

\begin{itemize}
\item
  First, the graph contains an explicit organizational layer through
  domains and ownership.
\item
  Second, admission into the triple set is mediated through governance
  rather than direct insertion.
\end{itemize}

In implementation, this is realized through a two-phase protocol:
propose a triple, obtain a governance decision, and then commit that
decision. This changes the graph from a passive fact store into a
managed memory structure.

The ownership function is written as φ: E → 2\^{}D, rather than φ: E →
D, because one entity may legitimately belong to more than one domain.
In a scientific corpus, the same entity can participate in multiple
conceptual regions of the graph: a method may belong to a methods domain
and also sit inside a task-specific domain; a biomarker may belong to
both an inflammation domain and a cardiometabolic domain. Using 2\^{}D
means φ(e) returns a set of owners, not a single owner. This is what
makes cross-domain governance possible. Once entities have domain
assignments, a triple can be classified as single-owner, cross-domain,
or unowned depending on the ownership of its endpoints.

Operationally, the data structure consists of more than entities and
triples. A governed triple also carries an ownership assignment, a
governance action, a rationale, and an audit record. In the
implementation, every admitted triple passes through
propose\_triple(...), produces a GovernanceDecision, and is then
committed through commit\_decision(...). This is the concrete mechanism
that instantiates $\gamma$. The graph therefore stores not only what is known,
but also how that knowledge entered memory and who was responsible for
it.

The main research idea is therefore not multi-agent extraction in
isolation, but the construction of a domain-governed organizational
memory in which knowledge is admitted, revised, and used through domain
ownership.

\textbf{System Overview}

The system has three layers.

\begin{itemize}
\item
  The first is the creation layer, which turns documents into entities
  and triples.
\item
  The second is the governance layer, which organizes knowledge into
  domains and controls admission into memory.
\item
  The third is the application layer, which uses the governed graph for
  downstream tasks such as question answering and incremental
  enrichment.
\end{itemize}

\textbf{Knowledge Graph Creation Workflow}

The system receives a corpus of documents. A top-level orchestrator
agent manages the end-to-end creation process over the full corpus and
over each document instance.

The pipeline is mostly sequential. Each specialized agent transforms or
validates the current representation and passes it to the next stage.
\st{This is intentionally simple.} The system only activates
deliberation when an agent signals uncertainty or conflict.

The workflow is as follows.


\begin{algorithm}[t]
\caption{Governed KG Creation}
\label{alg:gkg-creation}
\begin{algorithmic}[1]
\REQUIRE Corpus $\mathcal{C}$, Governed KG $G=(E,T,D,\phi,\gamma)$
\FOR{each document $x \in \mathcal{C}$}
    \STATE $\text{segments} \leftarrow \textsc{DocumentProcessor}(x)$
    \STATE $\text{schema} \leftarrow \textsc{DomainClassifier}(\text{segments})$
    \STATE $D_{\text{boot}} \leftarrow \textsc{BootstrapDomains}(\text{schema})$
    \STATE $\text{entities} \leftarrow \textsc{EntityExtractor}(\text{segments}, \text{schema})$
    \STATE $\phi \leftarrow \textsc{AssignEntitiesToDomains}(\text{entities}, D_{\text{boot}})$
    \STATE $\text{triples} \leftarrow \textsc{RelationExtractor}(\text{segments}, \text{entities}, \text{schema})$
    \STATE $\text{linked} \leftarrow \textsc{EvidenceLinker}(\text{triples}, \text{segments})$
    \STATE $\text{verified} \leftarrow \textsc{VerificationAgent}(\text{entities}, \text{linked}, x)$
    \FOR{each triple $t$ in $\text{verified}$}
        \STATE $\delta \leftarrow \textsc{propose\_triple}(t)$
        \STATE $\textsc{commit\_decision}(\delta)$
    \ENDFOR
\ENDFOR
\RETURN Governed KG $G$
\end{algorithmic}
\end{algorithm}

\begin{enumerate}
\def\labelenumi{\arabic{enumi}.}
\item
  The DocumentProcessor segments each document into coherent chunks.
\item
  The DomainClassifier derives a working schema for the document. In
  benchmark mode it uses a fixed schema; in open-world mode it discovers
  entity and relation types from the content. From this schema, the
  system bootstraps preliminary domains before any triples are admitted.
  This is the first point at which governance becomes active during
  creation.
\item
  The EntityExtractor identifies entity mentions, canonical labels, and
  candidate types. The shortest-meaningful-noun-phrase rule is used only
  to normalize the entity span itself, so that entities are anchored to
  compact canonical labels rather than overly long surface forms.
  Relation extraction is not limited to noun pairs: downstream agents
  still use the full sentence context, including predicates, modifiers,
  and supporting evidence, to determine whether a typed relation should
  be proposed between the extracted entities. The extractor is also
  instructed to preserve short scientific terms, abbreviations, tool
  names, and task phrases rather than discarding them as generic text.
\item
  Once entities exist, they are assigned to one or more domains. This
  partially instantiates the ownership function ($\phi$), allowing the system
  to know who owns a proposed fact during creation.
\item
  The RelationExtractor identifies relations between current entities
  and produces candidate triples. In fixed-schema benchmark mode, the
  extractor is constrained to the benchmark label set; in open-world
  mode, it can propose a new relation structure.
\item
  The EvidenceLinker grounds each candidate triple in the source
  document by retrieving explicit supporting text spans and source
  references. It is instructed to quote or localize evidence directly
  from the document rather than generate free-form justifications, and
  these grounded spans are then inspected by the verification and
  governance stages before the triple can be admitted.
\item
  The VerificationAgent removes unsupported or low-quality extractions
  and outputs a verified set of triples suitable for admission review.
\item
  At that point, the KnowledgeOrganizer and governance layer take over.
  Each verified triple is proposed to the governed graph rather than
  inserted directly. The system routes the triple to the relevant domain
  or domains, produces a governance decision, commits the decision, and
  appends that decision to the audit log.
\end{enumerate}

This is the step that distinguishes governed creation from ordinary
extraction.

There are several important operational details that the workflow above
relies on:

1. The corpus is processed incrementally. After each document, newly
accepted entities and triples are committed to the shared KG. Later
documents therefore see a richer memory than earlier documents.

2. Entity resolution happens against both the current document and the
pre-existing graph. A new mention is compared to entities already
extracted in the current document, entities already present in the KG,
and aliases stored in shared memory. This allows later documents to
attach to canonical entities created earlier in the corpus.

3. Domain bootstrapping begins before triple admission, but ownership
remains provisional as the graph grows. Early domains function as
lightweight organizational shells rather than fixed ontologies. New
facts are proposed from the source document first and only then
evaluated against the current graph, so existing triples do not
deterministically force new information into pre-existing relation
patterns. Instead, prior domain structure is used as review context: it
helps the system assess whether a candidate fact is consistent with
local ownership, nearby evidence, and previously admitted patterns,
while still allowing revision, cross-domain assignment, or rejection
when the new evidence does not fit the current graph. As a result, MaKG
differs from a traditional static KG by treating structure as governed
and revisable rather than fixed in advance.

4. Governance happens after extraction-time deliberation, not instead of
it. Deliberation asks whether an uncertain extraction should survive.
Governance asks how a surviving triple should enter memory.

What happens when a new document arrives can therefore be described
precisely. The system does not create a separate mini-KG and merge it
blindly. It segments the document, derives a schema, extracts entities,
resolves those entities against the existing graph and alias store,
assigns them to one or more domains, extracts candidate triples,
attaches evidence, filters low-quality proposals, and only then proposes
each surviving triple to the governed memory. The routing step uses the
ownership assignments of the subject and object to determine whether the
triple is single-owner, cross-domain, or effectively unowned. The
governance regime then determines whether the triple is automatically
accepted after routing, or whether it must be actively reviewed by the
responsible domain expert(s).

The post-routing cases should be stated explicitly:

- Single-owner triple: both endpoints resolve to the same governing
domain, so that domain becomes the primary owner.

- Cross-domain triple: the endpoints belong to different domains, so the
triple is treated as shared knowledge between owners.

- Unowned triple: no stable owner can be derived from the current org
chart, so the triple is escalated or logged accordingly.

The governance action itself can be one of five outcomes:

- approve: accept the triple as-is

- reject: do not admit it to memory

- revise: admit a normalized or corrected version of the triple

- escalate: mark it for explicit review because ownership or evidence is
insufficient

- auto\_approve: admit it automatically under a lightweight governance
regime while still recording the event

\textbf{Deliberation, Blackboard, and Voting}

The pipeline is sequential by default. Deliberation is an exceptional
path triggered when an agent reports low confidence. In the current
implementation, a threshold around 0.7 is used as the key decision
boundary.

If an agent's confidence is above threshold, the pipeline continues. If
confidence falls below threshold, the agent raises a hypothesis.

A hypothesis is a structured claim about an extraction unit that is not
trusted enough for automatic continuation. Typical hypothesis types are
an uncertain entity, an uncertain triple, or an uncertain relation type.
For example, from the sentence ``We present Minimum Bayes-Risk (MBR)
decoding for statistical machine translation,'' the system may first
propose the candidate triple (Minimum Bayes-Risk, Used-for, statistical
machine translation). That candidate is plausible, but the subject span
is not yet stable enough for admission because the text refers more
precisely to Minimum Bayes-Risk decoding. Rather than storing the weakly
normalized triple or discarding it outright, MaKG keeps it in a
hypothesis state so that evidence linking and governance review can
refine it before admission. This intermediate state is necessary because
many low-confidence cases are not incorrect; they are only locally
ambiguous, and resolving them immediately preserves document context
without polluting the graph with unstable structure.

Each hypothesis includes the proposed content, the posting agent, the
initial confidence, the supporting evidence, and the current pipeline
context.

The blackboard is a shared forum inside the shared memory system. It is
not a chat interface; it is a structured coordination object.

An agent posts a hypothesis to the blackboard, and the deliberation
coordinator asks a set of evaluator agents to vote on it.

Each voter is not necessarily a topical expert, but an expert on a
criterion.

The EntityExtractor evaluates whether the entity hypothesis is
plausible.\\
The RelationExtractor evaluates whether the relation structure is
coherent.\\
The EvidenceLinker evaluates whether the text supports the proposal.

This means the value of a vote comes from specialized perspective, not
topical ownership. Deliberation here is about extraction reliability.

Domain governance happens later, when the triple enters the governed KG.

Each voting agent returns a vote with rationale. The coordinator
aggregates the votes and resolves the hypothesis into a keep or no-keep
decision. If votes conflict, the coordinator can trigger a short
structured debate.

The final decision and rationale are written back to shared memory,
creating an audit trail that is used both operationally and
analytically. Operationally, the system reuses these records as
domain-local memory: prior approvals, rejections, and revisions help
later reviewers judge whether a new candidate is consistent with
previously admitted patterns or should be revised. Analytically, the
same trace makes it possible to inspect why a fact was admitted,
rejected, or modified, which supports ablation, error analysis, and
human validation.

The following example illustrates how hypothesis resolution works during
construction.\\
Consider the sentence: ``We present Minimum Bayes-Risk (MBR) decoding
for statistical machine translation.'' During extraction, the system
proposes the candidate triple (Minimum Bayes-Risk, Used-for, statistical
machine translation) with confidence 0.62, below the 0.70 admission
threshold. The candidate is promising, but not yet reliable enough to
store because the subject span may be under-specified: the text refers
to the specific method Minimum Bayes-Risk decoding, not just Minimum
Bayes-Risk in isolation. The candidate is therefore placed into a shared
hypothesis workspace for immediate review. The EntityExtractor votes
Revise (0.78): ``The method name should include decoding to match the
source span.'' The RelationExtractor votes Accept (0.82): ``The phrase
for statistical machine translation strongly supports a Used-for
relation.'' The EvidenceLinker votes Strong Accept (0.91): ``The source
sentence explicitly states the method-task relationship.'' The combined
outcome is to revise the subject and admit the corrected triple (Minimum
Bayes-Risk decoding, Used-for, statistical machine translation). If the
evidence had not supported the revision, the candidate would have been
rejected rather than inserted as a weakly normalized fact.

\textbf{Shared Memory}

MaKG uses shared memory to coordinate extraction within a document and
to maintain consistency across documents. These functions are separated
because the system must handle both short-lived local uncertainty and
longer-lived accumulated knowledge. During extraction, agents need
access to transient states such as pending hypotheses, intermediate
votes, and active evidence spans. Across documents, the system needs
persistent state such as admitted entities, aliases, domain ownership
patterns, and prior governance decisions.

In practice, the memory system has four regions.\\
\textbf{Working memory} stores the current pipeline state for the active
document, including candidate entities, candidate triples, pending
hypotheses, and vote requests. For example, if the system proposes
(Minimum Bayes-Risk, Used-for, statistical machine translation) but is
unsure whether the subject should be normalized more precisely as
Minimum Bayes-Risk decoding, that unresolved candidate lives in working
memory until review is complete. \textbf{Semantic memory} stores
accepted graph structure and canonicalized entities. Once the candidate
is revised and admitted, the canonical entity Minimum Bayes-Risk
decoding, its aliases such as MBR decoding, and the admitted triple are
written to semantic memory for reuse in later extraction, governance,
and QA.\\
\textbf{Procedural memory} stores reusable review patterns derived from
prior decisions, such as relation signatures that were previously
accepted, rejected, or revised. For example, if earlier decisions show
that ``Used-for'' should connect a method to a task rather than to a
vague concept, that pattern can inform later governance review.\\
\textbf{Episodic memory} stores per-document processing traces, such as
what was proposed, revised, or rejected in a specific document. This is
primarily useful for audit, debugging, and retrospective analysis rather
than as the main driver of graph construction.

These memory types are used at different timescales. Working memory is
short-lived and is reused continuously during a single document pass.
Semantic and procedural memory persist across documents and are
consulted repeatedly during governance and downstream QA. Episodic
memory is consulted less frequently, mainly when the system or a human
needs to inspect how a prior decision was made. This separation is
important because MaKG is not only extracting facts; it is building a
governed memory that must distinguish temporary uncertainty from
persistent knowledge.

\textbf{Output Schema}

The output of the creation pipeline is a JSON serialization of the
governed graph.

At minimum, it contains the knowledge graph, the org chart, the
governance mode, and the audit log.

Each entity stores an ID, labels, a type, and metadata (such as source
document and confidence).

Each triple stores a subject, relation, object, confidence, source, and
metadata.

The org chart is the organizational view of the graph. Each domain
includes:

\begin{itemize}
\item
  A domain identifier
\item
  A label and description
\item
  Seed entity and relation types
\item
  Current entity membership
\item
  A domain-specific subgraph
\item
  Cross-domain relations
\end{itemize}

The org chart summary is a compressed textual view listing domain names,
descriptions, relation signatures, and entity coverage.

\textbf{Applications Built on Top of the Governed KG}

The governed KG is the core object in MaKG, and downstream applications
operate over this governed memory rather than over raw extracted text
alone.

One application is question answering. A QA orchestrator first reads the
organizational structure of the graph, including the domain summary and
ownership layout, and uses that structure to route the question to the
most relevant domain experts. Each domain expert answers from its local
subgraph, grounded in the entities, relations, and nearby evidence it
owns. A synthesis stage then combines these domain-local responses into
a final answer. In this way, QA is not treated as flat retrieval over
the entire graph, but as a governed routing problem over domain-owned
knowledge.

A second application is incremental enrichment and self-organization. As
new documents arrive, candidate entities and triples are extracted and
compared against the existing graph. Rather than inserting them
directly, the system routes them through the same governance process
used during initial construction: ownership is assigned, evidence is
checked, prior domain memory is consulted, and the candidate is
approved, revised, or rejected before admission. This allows the graph
to grow over time while reorganizing itself around new evidence, new
entities, and new cross-domain links. In that sense, the graph is not
only incrementally enriched; it is also self-organizing under governed
update.

\subsection{Experimental Setup}\label{experimental-setup}

We evaluate MaKG along four empirical axes.

First, we evaluate KG construction quality on the SciERC test split in
fixed-schema mode. SciERC provides gold entity spans, entity types, and
typed relations over scientific abstracts. In this setting, the
extraction schema is fixed to the SciERC label set so that all systems
are evaluated against the same target ontology. This gives hard
benchmark metrics: entity F1, strict triple F1, fuzzy/mapped triple F1,
and triple hallucination. Strict triple F1 requires exact agreement on
subject, relation, and object. Fuzzy/mapped triple F1 is more forgiving
about surface normalization and relation-label canonicalization; this is
important because many apparent errors are canonicalization mismatches
rather than genuinely unsupported facts.

Second, we run governance ablations over the same 50-document GPT-5
candidate proposal stream. These ablations ask whether MaKG is better
because of governance itself or merely because some other part of the
pipeline filters triples. The ablations remove different pieces of the
governance stack: source filtering only, deterministic domain rules, a
centralized global reviewer, and finally full domain-memory governance.

Third, we prepare a human triple review study. This is necessary because
SciERC is a closed benchmark annotation, not an exhaustive list of every
valid scientific fact in an abstract. A triple can be supported by the
source text and useful for the graph while still being counted as a
false positive under SciERC if it is absent from the gold annotation.
The human review study samples governed-only, flat-only, shared, and
revised triples and asks annotators to judge support, relation
correctness, endpoint correctness, graph usefulness, and governance
decision quality.

Fourth, we evaluate downstream QA on MuSiQue, a multi-hop QA dataset.
MuSiQue is useful here because it tests whether the graph can support
chaining across documents rather than only retrieving one local
sentence. Unlike the SciERC extraction benchmark, this setting does not
provide a fixed target ontology, so MaKG is run in an open-world mode in
which entity and relation structure is induced from the corpus rather
than constrained to a predefined schema. We run a 50-question MuSiQue
slice with supporting-only contexts and compare MaKG domain-routed QA
against graph-based baselines, especially official Microsoft GraphRAG.
This tests whether the governed graph can be used as a stronger QA
substrate than a standard graph-based retrieval system.

\textbf{Results}

The completed results support four claims.

\begin{enumerate}
\def\labelenumi{\arabic{enumi}.}
\item
  Governance improves the admitted graph relative to flat insertion.
\item
  The improvement is not explained by source filtering or deterministic
  domain rules alone.
\item
  Domain experts become strongest when they carry domain-local memory.
\item
  The governed graph is useful for downstream QA relative to another
  graph-based system.
\end{enumerate}

\subsection{1. KG Construction Quality}\label{kg-construction-quality}

The first result compares direct flat KG insertion against two governed
variants of MaKG on the same 50 SciERC test documents using GPT-5. All
three systems share the same extractor, model, fixed SciERC schema, and
document set. The only difference is how candidate triples are admitted
into the stored graph.

The Flat KG baseline directly inserts extracted triples into a single
global graph. It has no ownership layer, no governance decision, no
revision step, and no domain-local memory.

MaKG basic governance routes each candidate triple to the owning domain
and applies governed admission before commit. The reviewer sees the
candidate triple, its source evidence, the ownership assignment derived
from the subject and object, and the endpoint neighborhood in the
current graph. The reviewer can approve, reject, or revise the triple.
This tests whether governance already improves graph quality even before
adding richer domain memory.

MaKG domain-memory governance is the full system. It keeps the same
governed routing and evidence-based admission, but the domain expert
also carries local memory: previously accepted and rejected relation
patterns, recent revised examples, endpoint membership in the owning
domain, and nearby admitted triples in the domain subgraph. This is the
strongest version of the claim in the paper, because it turns governance
from a generic review step into a domain-owned memory process.

\textbf{Table 1- GPT-5 KG Construction Quality On 50 SciERC Test
Documents}

{\def\LTcaptype{none} % do not increment counter
\begin{longtable}[]{@{}llllll@{}}
\toprule\noalign{}
\textbf{System} & \textbf{Entities} & \textbf{Triples} & \textbf{Strict
Triple F1} & \textbf{Fuzzy/Mapped Triple F1} & \textbf{TP / FP / FN} \\
GPT-5 Flat KG & 703 & 1366 & 0.116 & 0.204 & 100 / 1266 / 251 \\
GPT-5 MaKG basic governance & 687 & 750 & 0.144 & 0.305 & 79 / 671 /
272 \\
GPT-5 MaKG domain-memory governance & 687 & 472 & 0.173 & 0.333 & 71 /
401 / 280 \\
\midrule\noalign{}
\endhead
\bottomrule\noalign{}
\endlastfoot
\end{longtable}
}

The flat KG has the highest raw triple count, but most of those triples
do not match the SciERC gold graph. It produces 1366 triples, of which
1266 are counted as false positives under strict SciERC matching. MaKG
basic governance is substantially more selective: it reduces the
admitted graph to 750 triples, improves strict triple F1 from 0.116 to
0.144, improves fuzzy/mapped triple F1 from 0.204 to 0.305. This already
shows that governance changes what enters memory rather than acting as a
cosmetic layer on top of the same graph.

The full domain-memory version is stronger again. It admits 472 triples,
improves strict triple F1 to 0.173, and improves fuzzy/mapped triple F1
to 0.333. In other words, adding domain-local memory on top of governed
routing gives another clear lift over basic governance. This is the key
result for the paper: governance is most effective when domains are not
only routing destinations but active memory-holding experts.

The absolute F1 scores are still modest because SciERC is a strict
closed-label benchmark and does not treat all source-supported extracted
facts as true positives. For that reason, the paper should not claim
that extraction is solved. Given the same scientific extraction setting,
governed admission produces a cleaner admitted graph than flat
insertion, and domain-memory governance is the strongest admission
policy tested. Basic governance is worth keeping in the table because it
makes the mechanism legible: routing plus review already helps, and
routing plus domain-local memory helps more.

\subsection{2. Governance Ablation}\label{governance-ablation}

The purpose of this experiment is to isolate the value of governance
itself. All conditions start from the same 50-document GPT-5 candidate
proposal stream recorded during MaKG construction. This means the
extractor, documents, schema, and candidate facts are held fixed; only
the admission policy changes. The experiment therefore asks a focused
question: once candidate triples have been extracted, what is the best
way to decide which triples enter memory?

We compare three conditions.

The Flat KG baseline has no governance at all: extracted triples are
inserted directly into a single global graph.

The Global reviewer baseline applies a single centralized LLM reviewer
to the candidate triples. That reviewer only decides whether to approve
or reject each triple based on the triple and its evidence. It has no
notion of domain ownership, no domain-local memory, and no
subgraph-specific context. This is the cleanest centralized-review
baseline.

MaKG applies the full governed admission protocol. Each candidate triple
is routed to the owning domain, reviewed by the responsible domain
expert, and evaluated using not only the source evidence but also
domain-local memory: endpoint membership, previously accepted and
rejected relation patterns, and nearby admitted triples in the domain
subgraph. Only approved or revised triples are committed to memory.

This is the key distinction. The global reviewer is a generic filter
over triples. MaKG is a governed memory system: the review decision
depends on who owns the fact, what that domain already knows, and how
the new triple fits into the existing local subgraph.

\textbf{Table 2- Governance Ablation On 50 SciERC Test Documents}

{\def\LTcaptype{none} % do not increment counter
\begin{longtable}[]{@{}llllll@{}}
\toprule\noalign{}
\textbf{Condition} & \textbf{Admission policy} & \textbf{Triples} &
\textbf{Strict F1} & \textbf{Fuzzy/Mapped F1} & \textbf{TP / FP / FN} \\
Flat KG & direct insertion, no governance & 1366 & 0.116 & 0.204 & 100 /
1266 / 251 \\
Global reviewer & single centralized approve/reject reviewer & 674 &
0.146 & 0.304 & 75 / 599 / 276 \\
MaKG & domain-routed governance with domain-local memory & 472 & 0.173 &
0.333 & 71 / 401 / 280 \\
\midrule\noalign{}
\endhead
\bottomrule\noalign{}
\endlastfoot
\end{longtable}
}

The results show a clear progression. Moving from flat insertion to a
single global reviewer already improves the admitted graph
substantially: strict F1 rises from 0.116 to 0.146, fuzzy/mapped F1
rises from 0.204 to 0.304. This confirms that some form of review is
better than no review.

However, MaKG performs best overall. Compared with the global reviewer,
MaKG further improves strict F1 from 0.146 to 0.173, improves
fuzzy/mapped F1 from 0.304 to 0.333. The graph admitted by MaKG is
therefore not only smaller, but cleaner and more precise.

The interpretation is important. This is not evidence that ``more
prompting'' alone solves the problem. The global reviewer is already a
strong prompted filter. What MaKG adds is structured governance:
ownership, routing, local memory, and subgraph-aware review. Domain
experts do not review triples in isolation. They evaluate them relative
to what their domain already contains and what relation patterns have
previously been accepted or rejected. That additional organizational
context is what pushes MaKG past a single global reviewer.

This is the main ablation result- governance is not merely a binary
filter layered on top of extraction. A centralized reviewer helps, but a
governed, domain-owned memory helps more.

\subsection{3. Human Triple Review}\label{human-triple-review}

Tbd

\textbf{4. Downstream QA on MuSiQue}

To evaluate whether the governed graph is useful for downstream
applications, we test it on MuSiQue, a multi-hop question answering
benchmark designed so that answering a question requires combining facts
across multiple pieces of evidence rather than retrieving a single local
span. This makes it a good fit for our claim: the value of a governed
knowledge graph is not just that it stores facts, but that it organizes
them in a way that supports structured reasoning and routing. In
contrast to the SciERC experiments, this setting is open-world: MuSiQue
does not impose a fixed benchmark schema, so MaKG must organize
entities, relations, and domains directly from unconstrained text.

We use a 50-question slice from MuSiQue 2-hop. In this setting, each
question requires linking two supporting facts. For our experiment, the
supporting-only corpus contains 100 contexts drawn from 99 unique
titles. From this corpus, MaKG first constructs a governed graph in
open-world mode by extracting entities and relations from the passages,
organizing them into domain-owned subgraphs, and admitting facts through
its governance pipeline. We then build a QA system on top of this graph:
each question is decomposed and routed to the most relevant domain
experts, those experts answer from their local subgraphs and neighboring
evidence, and a final synthesis step combines their responses into a
single answer. We compare this domain-routed QA system against official
Microsoft GraphRAG, which is the cleanest graph-based baseline for the
paper. This comparison avoids confounding the result with plain
document-RAG reading-comprehension effects and instead asks the relevant
systems question: given a graph-oriented retrieval setup, does a
governed graph support better QA than a standard graph-based system?

In both systems, the model is GPT-5 and the evaluation uses the same 50
questions and the same gold answers. We report exact match (EM),
token-level F1, and answer rate. Exact match measures whether the
predicted answer string matches the gold answer after normalization.
Token F1 is more forgiving and captures partial overlap. Answer rate
measures how often the system returns an answer at all.

\textbf{Table 4- MuSiQue QA, 50-Question Slice}

{\def\LTcaptype{none} % do not increment counter
\begin{longtable}[]{@{}llllll@{}}
\toprule\noalign{}
\textbf{System} & \textbf{Setting} & \textbf{EM} & \textbf{Token F1} &
\textbf{Answer Rate} & \textbf{Interpretation} \\
GPT-5 GraphRAG & supporting-only & 0.580 & 0.709 & 1.00 & official
graph-aware baseline \\
GPT-5 MaKG domain QA & supporting-only & 0.660 & 0.775 & 1.00 & governed
graph with domain-routed QA \\
\midrule\noalign{}
\endhead
\bottomrule\noalign{}
\endlastfoot
\end{longtable}
}

MaKG domain QA outperforms official GraphRAG on both exact match and
token F1: 0.660 EM versus 0.580, and 0.775 token F1 versus 0.709. Both
systems answer every question, so the difference is not due to
abstention or failure to return outputs. The difference is in answer
quality.

This result supports the application-layer claim of the paper. GraphRAG
builds and retrieves from a graph-based representation, but it does not
preserve the same organizational structure as MaKG: ownership,
domain-local memory, governed revision history, and routing through
domain experts. MaKG uses the organization of the graph directly during
QA. Questions are routed through domain-owned subgraphs, domain experts
return subgraph-grounded answers, and the final response is synthesized
from these domain-local views. In effect, the graph is not only being
retrieved from; it is being used through its governance structure.

The result should still be interpreted carefully. This is a
supporting-only MuSiQue setup rather than a full open-retrieval
distractor setting, so the benchmark is not testing maximum retrieval
difficulty. What it does test is whether, once relevant information is
in the graph-building corpus, a governed graph supports better multi-hop
question answering than a standard graph-based retrieval system. On that
question, the answer is yes: the governed graph yields better QA
performance than GraphRAG.

\textbf{5. Token And Runtime Cost}

We also report computational cost because MaKG trades additional
governance calls for a cleaner and more auditable graph. The following
table uses exact API call and token counts from JSONL usage logs.

\textbf{Table 50 LLM Call And Token Cost}

{\def\LTcaptype{none} % do not increment counter
\begin{longtable}[]{@{}llllllll@{}}
\toprule\noalign{}
\textbf{Stage} & \textbf{Unit} & \textbf{LLM Calls} & \textbf{Prompt
Tokens} & \textbf{Completion Tokens} & \textbf{Total Tokens} &
\textbf{Calls / Unit} & \textbf{Tokens / Unit} \\
SciERC extraction build & 50 docs & 1169 & 4,042,896 & 814,849 &
4,857,745 & 23.38 / doc & 97,155 / doc \\
Domain-memory governance review & 1304 candidate triples & 49 & 754,434
& 83,183 & 837,617 & 0.038 / candidate & 642 / candidate \\
Full best extraction pipeline & 50 docs & 1218 & 4,797,330 & 898,032 &
5,695,362 & 24.36 / doc & 113,907 / doc \\
Global reviewer, matched & 1304 candidates & 49 & 97,829 & 70,751 &
168,580 & 0.038 / candidate & 129 / candidate \\
Global reviewer, minimal & 1304 candidates & 66 & 107,287 & 53,103 &
160,390 & 0.051 / candidate & 123 / candidate \\
MuSiQue KG build & 99 supporting titles & 328 & 3,348,277 & 430,130 &
3,778,407 & 3.31 / title & 38,166 / title \\
MuSiQue KG-QA ablation & 200 KG answers & 1371 & 3,349,454 & 2,080,567 &
5,430,021 & 6.86 / answer & 27,150 / answer \\
\midrule\noalign{}
\endhead
\bottomrule\noalign{}
\endlastfoot
\end{longtable}
}

Runtime was also material. The GPT-5 SciERC governed extraction build
took 9962 seconds, or 2.77 hours, for 50 documents. The flat GPT-5
SciERC build took 3891 seconds, or 1.08 hours. The 10-document
evidence-only ablation took 821 seconds. The current usage logs do not
separate every internal extraction sub-agent into distinct cost buckets.
Entity extraction, relation extraction, evidence linking, and any
deliberation triggered inside extraction are included in the
extraction-build usage log. Governance review is separately logged
because it was replayed as its own domain-memory review stage.

\textbf{Main Takeaway}

Flat insertion produces a larger graph, but it admits many noisy
triples. MaKG improves the admitted graph by adding governance, and the
best result appears when domain experts carry local memory rather than
acting as generic reviewers. The ablation shows that source filtering
alone, deterministic ownership alone, and a centralized global reviewer
do not fully explain the result. The full domain-memory governance
condition is the strongest construction system tested: it has the best
strict F1, best fuzzy/mapped F1, and lowest hallucination. Downstream,
MaKG outperforms GraphRAG on MuSiQue graph-based QA, showing that the
governed graph can serve as a useful application substrate rather than
only a passive storage artifact. The human review study is the final
validation layer, designed to show whether triples counted as false
positives by SciERC are in fact source-supported and useful for the
governed graph.
